\section{Preliminaries}
\label{sec:prelim}

We address flow-level intrusion detection on IoT gateways. A flow
$\mathbf{x} \in \mathbb{R}^{d}$ is a fixed-length vector of $d$ numeric
descriptors derived from a network connection. A detector $f$ maps
$\mathbf{x}$ to one of $K$ known attack-family labels in
$\mathcal{K} = \{c_1, \dots, c_K\}$. The closed-set assumption is standard
but brittle, and Section~\ref{subsec:zeroday} formalises a retrieval-confidence
protocol that relaxes it for zero-day traffic.

Retrieval-augmented generation (RAG) conditions a language model on content
retrieved from an external store rather than on parametric memory alone. The
canonical pattern embeds a corpus into a dense vector index and retrieves
the top-$k$ neighbours of a query by cosine similarity as in-context
evidence. We specialise this pattern to policy synthesis. The corpus is the
training split of an IoT dataset, each flow is embedded under its class
label, and the model emits a structured policy rather than free text. In
the SMCC view of this special issue, the vector store and the language
model together act as the \emph{memory} stage. They compress the training
distribution into a retrieval index and a prompt context that are touched
once at offline policy-generation time.

A threshold rule is a tuple
\texttt{(feature, op, value, target\_class)} that fires when the named
feature of $\mathbf{x}$ satisfies the comparison. A policy is a set of such
rules, and the detector evaluates every rule against $\mathbf{x}$ in a
vectorised comparison and assigns the majority-vote class. Threshold rules
are the same primitive that IoT gateways already execute as access-control
list (ACL) firewall entries, so a generated policy deploys without new
runtime infrastructure. In the SMCC view, the rule evaluator is the
\emph{computation} stage. It runs at line rate, carries no model weights,
and issues no language-model call per flow.
Section~\ref{sec:architecture} composes these two primitives into the
end-to-end system, and Section~\ref{sec:framework} defines the four
protocols that exercise it.
